% Standalone note. Paste into the weekly report, or \input it, the same way
% thesis/week4/locked_tortuosity_metrics.tex is used.
%
% READING DISCIPLINE. No new reading was done for this note. Every number, citation and evidence
% tag below is carried over unchanged from two reports already written in this repo:
%   - reports/Coronary bifurcation diameter descriptors.md (the locked descriptor panel, its math,
%     its ImageCAS-X reliability numbers, and its clinical-evidence sourcing)
%   - reports/Bifurcation calibre in both models.md (how the panel is consumed by Model 1 and
%     Model 2)
% and their backing research_notes/ directories. Evidence tags follow those reports' own
% convention: [FT] full text read, [ABS] abstract only, [SNIP] search summary only, [DERIVED]
% worked out analytically, [SIM] phantom or Monte Carlo run, [MEASURED] run on ImageCAS-X data.
% Citations to external sources are given as direct links rather than \cite{} against
% thesis/refs.bib, because most of them have not yet been added there as verified bib entries;
% three that were already locked for the tortuosity note (taylor2024murray, asakura1990flow,
% chatzizisis2007ess) are cited normally, since they already exist as verified entries.
%
% ORDERING. Descriptors below are ranked by (1) how comparable the measurement is across trees at
% the same anatomical bifurcation (reliability and how often the junction exists at all), then (2)
% how plausible, not proven, a link to future coronary artery disease (CAD) is. This follows Kit's
% stated priority, "the main goal is the H\O{} dataset and matching the shape descriptors with
% outcome" and "I don't want you to get bogged down in the details of representation learning for
% now" (supervisor correspondence, from\_superviser.tex), and repeats this project's own precedent
% from locking the tortuosity panel: mean absolute curvature was locked as primary because it was
% the measure that survived a real evidence test ($R^2=0.112$), not the arc/chord tortuosity index
% the literature had focused on, which tested at $R^2\approx 0$ on the same data. Reliability first,
% literature enthusiasm second, applied again here.

\documentclass[11pt]{article}
\usepackage[margin=1in]{geometry}
\usepackage{amsmath}
\usepackage{textcomp}
\usepackage[colorlinks=true,linkcolor=blue,citecolor=blue,urlcolor=blue]{hyperref}
\usepackage{booktabs}
\usepackage{array}
\usepackage{microtype}

\title{Locked bifurcation calibre descriptors: clinical evidence, the math, and how they enter Model 1 and Model 2}
\author{}
\date{}

\begin{document}
\maketitle
\vspace{-2em}

\section*{The locked bifurcation calibre panel, and why each piece is in it}

A bifurcation is where one vessel splits into two, and it is the other canonical site, alongside a
sharp bend, where plaque tends to start: tracing flow through human coronary trees after death,
Asakura and Karino found almost every plaque at the inner wall of a bend or the outer wall just past
a bifurcation (\href{https://doi.org/10.1161/01.RES.66.4.1045}{Asakura and Karino 1990}), and
chronically low wall shear stress (WSS) is the mechanism already used to motivate the locked
tortuosity panel (\href{https://doi.org/10.1016/j.jacc.2007.02.059}{Chatzizisis et al.\ 2007}). This note extracts
the analogous locked panel for bifurcation \textbf{calibre}, meaning vessel radius, rather than
bending. It covers five descriptors, the covariates carried alongside them, two descriptors that were
considered and left out, and, new relative to the tortuosity note, exactly how each piece enters
this thesis's two comparison models.

\subsection*{Model 1 and Model 2, and the two zones every measurement avoids}

\textbf{Model 1} is a logistic regression on hand-built scalar descriptors (plus age, sex and
covariates); \textbf{Model 2} is a small network that reads a per-point profile along the vessel
directly, with the same covariates late-fused in. The whole point of building a fair comparison is
that neither model should see information the other is denied.

Every measurement below is read from the \textbf{EDT radius}: at each point on the vessel's
centerline, the distance from that point to the nearest vessel-wall voxel, which is the standard
inscribed-sphere estimate of the local vessel radius. Right at a bifurcation this measurement is
misleading. Picture the vessel as a garden hose briefly flaring into a Y-joint before narrowing back
into two branches: the EDT algorithm works by finding the largest circle that fits inside the vessel
at each point, and near the joint that circle can keep touching the far wall of the \emph{combined}
Y-shape, so it reads close to the parent's radius for a few millimetres in every direction, whether
or not the vessel actually widens there at all. Phantom testing confirms this is mostly a
computational artefact, not anatomy: on a bifurcation phantom built with no real flare, the ratio of
the junction reading to the true parent radius still sits at 0.93 to 0.99 [SIM]. Every
descriptor below is therefore read from a window that starts one \textbf{junction-sphere radius}
($r_J$, the EDT value at the branch point itself) away from the branch point, so the artefact zone is
skipped before any radius is averaged.

\subsection*{1. The left main (LM) Finet ratio, the primary measure}

The left main splits into the left anterior descending (LAD) and left circumflex (LCx). Writing
$r_0$ for the median EDT radius on the parent (LM) over a 3\,mm window starting $1{\cdot}r_J$
upstream of the split, and $r_1, r_2$ for the same measurement on the two daughters starting
$1{\cdot}r_J$ downstream, the first locked measure is
\begin{equation}
  F = \frac{r_0}{r_1 + r_2} .
  \label{eq:finet}
\end{equation}
Finet and colleagues found, across 173 bifurcations in 47 patients with normal coronary
angiograms, that the parent radius tracks a fixed fraction of the sum of the daughters,
$r_0 \approx 0.678\,(r_1+r_2)$
(\href{https://eurointervention.pcronline.com/article/fractal-geometry-of-arterial-coronary-bifurcations-a-quantitative-coronary-angiography-and-intravascular-ultrasound-analysis}{Finet
et al.\ 2008} [FT]). $F$ close to 0.678 means a bifurcation follows that classical proportion; a
lower $F$ means the parent reads narrow relative to its daughters. On 36 ImageCAS-X test-split
reference cases the median is $F=0.615$, and on 100 training-split cases $F=0.63$
[MEASURED], somewhat below Finet's 0.678 and below a CT atlas mean of $0.658\pm0.083$
(\href{https://doi.org/10.4244/EIJV12I7A139}{Medrano-Gracia et al.\ 2016} [FT]); the diameter report's error analysis shows this gap is largely a
measurement-convention effect of the EDT estimator, not necessarily anatomy.

No study has tested $F$ itself against coronary outcome. The closest real evidence is retinal:
individual-participant data on 22{,}159 people free of coronary disease found narrower retinal
arterioles carried a hazard ratio (HR) of 1.17 per 20\,\textmu m in women and 1.02 in men for later
coronary heart disease
(\href{https://pubmed.ncbi.nlm.nih.gov/19755365/}{McGeechan et al.\ 2009} [ABS]), and two further
cohorts gave HR 1.06 to 1.20 per standard deviation (SD) narrower
(\href{https://pubmed.ncbi.nlm.nih.gov/27682886/}{Seidelmann et al.\ 2016} [ABS];
\href{https://doi.org/10.1038/s41551-020-00626-4}{Cheung et al.\ 2021} [ABS]). This is borrowed from
a different vascular bed and a different disease pathway, not a coronary bifurcation precedent, and
should be read as a rough prior on effect size, not evidence for $F$ specifically.

\emph{Comparability.} Every coronary tree has exactly one LM bifurcation, so $F$ is measurable in
essentially every case, and it is the most reliable descriptor in the panel: intraclass correlation
coefficient (ICC) 0.85 for the EDT-versus-surface-distance method and 0.90 for the choice of
measurement window [MEASURED].

\subsection*{2. The daughter ratio, the best-evidenced measure mechanistically}

At any of the four named junctions this thesis measures (LM$\to$LAD/LCx, LAD$\to$D1,
LCx$\to$OM1, and the RCA crux $\to$PDA/PLB), the second locked measure is the smaller daughter
radius over the larger,
\begin{equation}
  \rho = \frac{\min(r_1, r_2)}{\max(r_1, r_2)} \in (0, 1] ,
  \label{eq:daughter-ratio}
\end{equation}
with $\rho=1$ a perfectly symmetric split and $\rho \to 0$ a highly asymmetric one. This is the
descriptor with the strongest mechanistic case in the whole panel. In 230 synthetic left-coronary
computational fluid dynamics (CFD) models, the area of vessel wall running under low WSS followed a
quadratic curve in the LAD/LCx diameter ratio with $R^2 = 0.91$ to $0.94$, and the low-shear area
formed preferentially in the \emph{smaller} branch
(\href{https://pmc.ncbi.nlm.nih.gov/articles/PMC11770140/}{Garcha and Grande Guti\'errez 2025}
[FT]). In 39 real human left coronary trees, diameter was the only piece of vessel geometry that
survived correction for multiple comparisons as a correlate of WSS
(\href{https://arxiv.org/abs/2502.06161}{Shen et al.\ 2025} [FT, preprint]). The physical intuition
is direct: an asymmetric split still has to carry the same total flow, so the smaller daughter
receives less flow than its own size alone would suggest is efficient, which is exactly the kind of
mismatch that produces chronically low WSS.

\emph{Comparability.} Best at the LM, ICC 0.86, degrading toward the distal junctions (ICC
approximately 0.63 to 0.79 at LAD-D1, LCx-OM1 and the crux) [MEASURED]. The LCx-OM1 junction is
missing in about 20 of every 100 trees, since OM1 is not always present, so this descriptor needs an
explicit absence flag rather than silent dropping.

\subsection*{3. The Huo-Kassab optimality-ratio deviation ($\Gamma_{HK}$), in place of the raw exponent}

The classical way to ask whether a bifurcation follows an optimal flow-diameter law is to solve for
the exponent $n$ in
\begin{equation}
  r_0^{\,n} = r_1^{\,n} + r_2^{\,n} .
  \label{eq:exponent}
\end{equation}
Murray's cube law sets $n=3$, derived from the assumption that the body minimises the combined cost
of pumping blood and maintaining vessel tissue, which is optimised when wall shear stress is
constant across the junction. In practice coronary trees run lower than this: pooling 1070 coronary
trees, the exponent's 95\,\% confidence interval sits at 2.24 to 2.54, centred on 2.39, distinctly
below 3 (\href{https://doi.org/10.1152/ajpheart.00142.2024}{Taylor et al.\ 2024}), closer to the Huo-Kassab value $n=7/3\approx2.33$, an extension of
Murray's law that accounts for how blood's effective viscosity itself depends on vessel diameter.

Equation~\ref{eq:exponent} is a poor number to compute directly. It has no closed-form solution, so
it must be found numerically, and it is \textbf{ill-conditioned}: a small radius-measurement error
gets amplified by a factor of roughly $n^2/H$, where $H$ measures how unevenly the flow is split
between the two daughters, and the equation has \emph{no solution at all} once a noisy daughter
radius reads as large as, or larger than, the parent's [DERIVED]. On the real LM data, exactly this
happened in 10 of 36 junctions, 28\,\% [MEASURED].

The fix borrowed here comes from the retinal-imaging literature, which solved the same
ill-conditioning problem with an \textbf{optimality ratio} rather than an exponent. For the cube law,
$\Gamma = \left(\tfrac{d_1^3+d_2^3}{2}\right)^{1/3}/d_0$, which equals $2^{-1/3}=0.794$ exactly at a
Murray optimum, for any daughter asymmetry, and which rose measurably under a drug that blocks
nitric-oxide synthase, showing it is sensitive to real endothelial function, not just noise
(\href{https://pmc.ncbi.nlm.nih.gov/articles/PMC2954284/}{Witt et al.\ 2010} [FT]). The coronary
version of this ratio, matching the coronary field's own $7/3$ exponent instead of the retinal cube,
is
\begin{equation}
  \Gamma_{HK} = \frac{1}{r_0}\left(\frac{r_1^{\,7/3} + r_2^{\,7/3}}{2}\right)^{3/7} ,
  \qquad\text{locked descriptor: } \Gamma_{HK} - 0.743 ,
  \label{eq:gamma-hk}
\end{equation}
where $2^{-3/7}\approx0.743$ is the value $\Gamma_{HK}$ takes exactly whenever the Huo-Kassab law
holds precisely, again for any asymmetry between the daughters [DERIVED here]. The locked descriptor
is the signed deviation $\Gamma_{HK}-0.743$: unlike $n$, it stays well-defined and numerically stable
under measurement noise, because, like $F$ and $\rho$ above, it is built entirely from ratios of
radii and inherits their well-behaved error algebra rather than the exponent's.

The clinical motivation is one step removed and worth stating carefully. In 253 patients imaged with
intravascular ultrasound, a higher value of a Murray-type ratio, in Schoenenberger and colleagues'
own convention describing how large the parent vessel reads relative to its daughters, was
associated with denser calcium at the bifurcation after adjustment for age, sex and standard risk
factors (\href{https://pubmed.ncbi.nlm.nih.gov/22261173/}{Schoenenberger et al.\ 2012} [ABS]). Their
ratio convention is the mirror image of $\Gamma_{HK}$ (parent over daughters, rather than daughters
over parent), so this finding is evidence that a law-deviation descriptor in this family is worth
including, not yet licence to predict which sign of $\Gamma_{HK}-0.743$ should carry the risk; that
sign should be checked, not assumed, once $\Gamma_{HK}$ is actually computed.

\emph{Comparability.} $\Gamma_{HK}$ is always mathematically defined, with none of $n$'s unsolvable
cases, but it is not yet implemented in this project's code and so has no measured ImageCAS-X ICC
yet; because it is built from the same two radii as the daughter ratio above, similar reliability is
expected but not confirmed.

\subsection*{4. Parent-to-daughter ratios, secondary and the most window-sensitive}

The two simple ratios that make up the Finet ratio's own numerator and denominator,
\begin{equation}
  \frac{r_{\mathrm{LAD}}}{r_{\mathrm{LM}}}, \qquad \frac{r_{\mathrm{LCx}}}{r_{\mathrm{LM}}} ,
  \label{eq:parent-daughter}
\end{equation}
read 0.884 and 0.774 on ImageCAS-X [MEASURED]. They carry no clinical evidence beyond what $F$ and
$\rho$ already capture, since they are literally $F$'s two components taken apart, and they are the
single most window-sensitive descriptor measured: moving the measurement window from the branch
point to this project's default offset shifted them by 0.10 to 0.12, about two-thirds of a
between-subject standard deviation [DERIVED, MEASURED]. They are kept as secondaries because
disaggregating $F$ into its two sides is occasionally informative, but they should never be reported
without the exact window definition stated alongside them.

\subsection*{5. Side-branch Finet ratio, exploratory}

The same formula as Equation~\ref{eq:finet}, applied at LAD-D1, LCx-OM1 and the RCA crux instead of
the LM. Reliability falls the further the junction sits from the LM: ICC roughly 0.79 at LAD-D1,
0.67 at LCx-OM1, and only 0.44 at the crux [MEASURED], the last of which is too unreliable to trust
as a modelled predictor at all. The crux additionally exists only in right-dominant trees, and OM1
is absent in about a fifth of trees. Kept as an exploratory, reported-only item, not a modelled
feature, at these junctions.

\subsection*{Covariates carried alongside the panel}

Three further quantities travel with the calibre panel as covariates rather than as calibre
descriptors in their own right: the fixed-tangent bifurcation angle, already this project's locked
primary geometric descriptor from the separate regional-geometry work; a binary ramus flag, since a
third LM daughter changes the bifurcation's whole geometry and cases with one should be reported as
their own stratum rather than pooled; and LM length (typically $10.5\pm5.3$\,mm), needed because a
short LM can clip the parent measurement window.

\subsection*{Considered, not locked}

The raw junction exponent $n$ (Equation~\ref{eq:exponent}) is kept as a descriptive number only,
never a modelled feature: it is unsolvable in 12 to 37\,\% of junctions and its ICC runs only 0.48 to
0.69 even where it can be computed [MEASURED]. An area-ratio analogue of Equation~\ref{eq:daughter-ratio},
$(r_1^2+r_2^2)/r_0^2$, was also considered and left out: its relative measurement error runs at
roughly twice that of a radius ratio for the same input noise, and it adds no information the radius
ratios above do not already carry.

\section*{How the panel enters Model 1 and Model 2}

The mechanism is the same one already locked for plain vessel diameter, repeated here at junction
scale: \textbf{both arms see the identical numbers, and neither arm sees information the other
lacks.} Concretely, at the LM this is an 8-field vector: the Finet ratio (primary), the daughter
ratio, the $\Gamma_{HK}$ deviation, the two parent-to-daughter ratios, the angle, the ramus flag, and
LM length. LAD-D1 and LCx-OM1 drop to 3 continuous fields (there is only one daughter to compute a
ratio against, and no LM-length equivalent) plus the angle, with LCx-OM1 adding an OM1-absence flag.
The crux keeps the same shape but stays reported-only, per its ICC 0.44.

\begin{itemize}
  \item \textbf{Model 1} receives every field as an extra regression term: continuous fields spline-
    transformed with a handful of pre-specified knots, the ramus flag as a stratum indicator rather
    than a pooled covariate.
  \item \textbf{Model 2} receives the identical numbers, late-fused into the network \emph{after} its
    pooled sequence embedding, in addition to the masked raw radius profile it already reads on each
    limb (the profile with the flare artefact zone removed, exactly as Model 1's windows are).
\end{itemize}
Giving Model 2 both the scalar and the raw profile is not an unfair advantage to the network: the
scalar is a fixed, deterministic function of the same profile the network already reads, so nothing
is handed to it that it could not in principle compute for itself, and nothing is withheld from
Model 1 that changes what it is allowed to see.

One place this differs from the whole-vessel design: for a full-length profile, an
``information-matched'' middle rung between Model 1 and Model 2, functional logistic regression on
the profile's leading principal components, was proposed elsewhere in this project's design. That
does not transfer safely to a 3\,mm junction window. At the 0.5\,mm resolution this project's data
is cached at, a 3\,mm window holds roughly six to seven points along the centerline, too few to fit a
stable functional model with more than one or two real modes of variation [DERIVED, not yet
measured]. The recommended substitute at junction scale is a short, fixed set of summary statistics
(the window median already used above, plus an endpoint-to-plateau ratio, since the diameter report
already measured exactly this quantity), not a functional-data method.

A learned alternative, a network that builds its own embedding of a junction from raw points rather
than being handed the explicit 8-field vector, is the design one supervisor (Bj\o rn) has asked for.
No coronary, cerebral or retinal precedent found in this project's research trains a junction
embedding that way; the one close match, a 2025 Circle-of-Willis pipeline, computes bifurcation
features with exactly the Finet and Huo-Kassab formulas already locked above, which supports the
descriptor \emph{content} chosen here but not a learned \emph{architecture}
(\href{https://arxiv.org/html/2510.13720v1}{Musio et al.\ 2025} [FT]). Consistent with Kit's stated
priority, the explicit vector above is the default design; a learned junction token stays a later,
optional ablation, not the starting point.

\section*{Priority order: comparable first, then a plausible CAD link}

This is a research question, so the CAD-link column below is a prior worth testing, not a
demonstrated association; no coronary bifurcation calibre descriptor has ever been tested against a
real coronary outcome.

\begin{table}[h]
\centering
\small
\begin{tabular}{@{}p{0.4cm}p{3.6cm}p{5.1cm}p{5.6cm}@{}}
\toprule
\# & Descriptor & Comparable across trees? & Plausible CAD link (a prior, not proof) \\
\midrule
1 & LM Finet ratio & Highest: present at every LM, ICC 0.85--0.90 & Weak--moderate, borrowed from retinal calibre (HR $\sim$1.05--1.20/SD); no coronary bifurcation-outcome study exists \\
2 & Daughter ratio & High at LM (ICC 0.86), lower distally (ICC $\sim$0.63--0.79); OM1 missing $\sim$20\% & Strongest in the panel: CFD $R^2$=0.91--0.94; only geometry surviving correction in a human WSS study \\
3 & $\Gamma_{HK}$ deviation & Always defined (no unsolvable cases); not yet measured on ImageCAS-X & Specific but thin: one IVUS study links a related ratio to bifurcation calcium; sign unresolved for this exact formula \\
4 & Parent-to-daughter ratios & Present at LM; most window-sensitive descriptor measured & None beyond what \#1 and \#2 already carry (they are \#1's own components) \\
5 & Side-branch Finet & Degrades away from LM: ICC 0.79 / 0.67 / 0.44; crux right-dominant only & None; untested transplant of the LM formula \\
--- & Exponent $n$, area ratio & \emph{Considered, not locked}: $n$ unsolvable in 12--37\,\% of junctions & Not usable as a modelled predictor regardless of any CAD evidence \\
\bottomrule
\end{tabular}
\end{table}

None of the five locked descriptors is a direct measurement of wall shear stress or of disease, and
a high or low score on one is a candidate risk marker, not proof of risk, exactly as already stated
for the tortuosity panel. The order above optimises for what this project can actually deliver first,
a reliable, reproducible measurement at a named anatomical location across the whole cohort, and
treats the coronary artery disease link as the open question the Herlev-\O{}sterbro outcome data will
have to answer, not one the literature has already settled.

\end{document}
